% GENERATED FILE. Edit canonical lessons, then run npm run generate:books.
% canonical-source-hash: fnv1a64:878a3cb17dea4194
% canonical-lessons: ES-C592-empresario, ES-C592-multinacional, ES-C592-industria, ES-C592-sede, ES-C592-sucursal

\chapter{La empresa}
\label{ch:es-la-empresa}
\hlchaptermodality{592}

\begin{chapteropening}
\textbf{By the end of this chapter:} \emph{I can say what kind of company someone works for and where its head office and branches are.}

Complete the last lesson of chapter 592: I can say what kind of company someone works for and where its head office and branches are.
\end{chapteropening}

\section[el empresario]{el empresario — the businessman, the employer}

\label{lesson:ES-C592-empresario}



\begin{quote}
Before the new one: say the Spanish for remote work, then the Spanish for the trade union.
\end{quote}

\subsection*{You'll want to know: el empresario}
\textbf{el empresario} — \textquotedblleft{}the businessman, the employer\textquotedblright{}.

\begin{quote}
Su padre es empresario.
\end{quote}

\begin{quote}
\emph{His father is a businessman.}
\end{quote}

A woman is \emph{la empresaria}. \emph{Los empresarios} can mean employers as a group.

\begin{cousinweb}
\emph{empresario} is built on \emph{empresa}, which was borrowed from Italian \emph{impresa} 'undertaking', built on Latin \emph{prehendere} 'to take hold of' (EtymDB).
\end{cousinweb}

\begin{grammarlens}[title={What you've built}]
One more word for telling, explaining and describing at B1.
\end{grammarlens}

\subsection*{Your turn}
\begin{itemize}
\raggedright
  \item \emph{Say it:} \emph{el empresario}
  \item \emph{Say it:} \emph{el empresario}, once more
  \item \emph{Say it:} a sentence of your own with \emph{el empresario}
  \item \emph{Recall:} say the Spanish for remote work, then the Spanish for the trade union, then say \emph{el empresario} again
\end{itemize}

\begin{tcolorbox}[breakable,colback=teal!4,colframe=teal!35!black,title={Before you move on}]
What does el empresario mean? (\textquotedblleft{}The businessman, the employer\textquotedblright{}.) Say it once more.
\end{tcolorbox}
\section[la multinacional]{la multinacional — the multinational company}

\label{lesson:ES-C592-multinacional}



\begin{quote}
Before the new one: say the Spanish for the trade union, then the Spanish for the businessman.
\end{quote}

\subsection*{You'll want to know: la multinacional}
\textbf{la multinacional} — \textquotedblleft{}the multinational company\textquotedblright{}.

\begin{quote}
Trabaja en una multinacional.
\end{quote}

\begin{quote}
\emph{She works for a multinational.}
\end{quote}

Also an adjective: \emph{una empresa multinacional}.

\begin{cousinweb}
\emph{multinacional} is built on \emph{nación}, a learned borrowing from Latin \emph{natio} 'birth, a people' (EtymDB), from \emph{nasci} 'to be born', with Latin \emph{multi-} 'many'.
\end{cousinweb}

\begin{grammarlens}[title={What you've built}]
One more word for telling, explaining and describing at B1.
\end{grammarlens}

\subsection*{Your turn}
\begin{itemize}
\raggedright
  \item \emph{Say it:} \emph{la multinacional}
  \item \emph{Say it:} \emph{la multinacional}, once more
  \item \emph{Say it:} a sentence of your own with \emph{la multinacional}
  \item \emph{Recall:} say the Spanish for the trade union, then the Spanish for the businessman, then say \emph{la multinacional} again
\end{itemize}

\begin{tcolorbox}[breakable,colback=teal!4,colframe=teal!35!black,title={Before you move on}]
What does la multinacional mean? (\textquotedblleft{}The multinational company\textquotedblright{}.) Say it once more.
\end{tcolorbox}
\section[la industria]{la industria — the industry}

\label{lesson:ES-C592-industria}



\begin{quote}
Before the new one: say the Spanish for the businessman, then the Spanish for the multinational company.
\end{quote}

\subsection*{You'll want to know: la industria}
\textbf{la industria} — \textquotedblleft{}the industry\textquotedblright{}.

\begin{quote}
La industria del coche es grande aquí.
\end{quote}

\begin{quote}
\emph{The car industry is big here.}
\end{quote}

\emph{Industrial} is the adjective.

\begin{cousinweb}
\emph{industria} is a learned borrowing from Latin \emph{industria} 'diligence, activity' (EtymDB).
\end{cousinweb}

\begin{grammarlens}[title={What you've built}]
One more word for telling, explaining and describing at B1.
\end{grammarlens}

\subsection*{Your turn}
\begin{itemize}
\raggedright
  \item \emph{Say it:} \emph{la industria}
  \item \emph{Say it:} \emph{la industria}, once more
  \item \emph{Say it:} a sentence of your own with \emph{la industria}
  \item \emph{Recall:} say the Spanish for the businessman, then the Spanish for the multinational company, then say \emph{la industria} again
\end{itemize}

\begin{tcolorbox}[breakable,colback=teal!4,colframe=teal!35!black,title={Before you move on}]
What does la industria mean? (\textquotedblleft{}The industry\textquotedblright{}.) Say it once more.
\end{tcolorbox}
\section[la sede]{la sede — the head office, the headquarters}

\label{lesson:ES-C592-sede}



\begin{quote}
Before the new one: say the Spanish for the multinational company, then the Spanish for the industry.
\end{quote}

\subsection*{You'll want to know: la sede}
\textbf{la sede} — \textquotedblleft{}the head office, the headquarters\textquotedblright{}.

\begin{quote}
La sede de la empresa está en Madrid.
\end{quote}

\begin{quote}
\emph{The company's head office is in Madrid.}
\end{quote}

\emph{La sede central} — the head office.

\begin{cousinweb}
\emph{sede} is a learned borrowing from Latin \emph{sedes} 'seat, home' (EtymDB), built on \emph{sedere} 'to sit'.
\end{cousinweb}

\begin{grammarlens}[title={What you've built}]
One more word for telling, explaining and describing at B1.
\end{grammarlens}

\subsection*{Your turn}
\begin{itemize}
\raggedright
  \item \emph{Say it:} \emph{la sede}
  \item \emph{Say it:} \emph{la sede}, once more
  \item \emph{Say it:} a sentence of your own with \emph{la sede}
  \item \emph{Recall:} say the Spanish for the multinational company, then the Spanish for the industry, then say \emph{la sede} again
\end{itemize}

\begin{tcolorbox}[breakable,colback=teal!4,colframe=teal!35!black,title={Before you move on}]
What does la sede mean? (\textquotedblleft{}The head office, the headquarters\textquotedblright{}.) Say it once more.
\end{tcolorbox}
\section[la sucursal]{la sucursal — the branch}

\label{lesson:ES-C592-sucursal}



\begin{quote}
Before the new one: say the Spanish for the industry, then the Spanish for the head office.
\end{quote}

\subsection*{You'll want to know: la sucursal}
\textbf{la sucursal} — \textquotedblleft{}the branch\textquotedblright{}.

\begin{quote}
El banco tiene una sucursal en mi barrio.
\end{quote}

\begin{quote}
\emph{The bank has a branch in my neighbourhood.}
\end{quote}

A branch of a bank or a shop.

\begin{cousinweb}
\emph{sucursal} was borrowed from French \emph{succursale}, which goes back to Medieval Latin \emph{succursus} 'help' (EtymDB links the Spanish word to \emph{succursus}), from \emph{succurrere} 'to run to help', built on \emph{currere} 'to run'.
\end{cousinweb}

\begin{grammarlens}[title={What you've built}]
One more word for telling, explaining and describing at B1.
\end{grammarlens}

\subsection*{Your turn}
\begin{itemize}
\raggedright
  \item \emph{Say it:} \emph{la sucursal}
  \item \emph{Say it:} \emph{la sucursal}, once more
  \item \emph{Say it:} a sentence of your own with \emph{la sucursal}
  \item \emph{Recall:} say the Spanish for the industry, then the Spanish for the head office, then say \emph{la sucursal} again
\end{itemize}

\begin{tcolorbox}[breakable,colback=teal!4,colframe=teal!35!black,title={Before you move on}]
What does la sucursal mean? (\textquotedblleft{}The branch\textquotedblright{}.) Say it once more.
\end{tcolorbox}
